\section{Further research questions}
\label{cont26:sec:further}

The following questions have concrete stopping conditions: an exact
identity, a stated uniform remainder, a certified classification, or
a counterexample would constitute progress. They are ordered by
their proximity to the present results, rather than by a claim that
the first questions are easy.

\subsection{Gaussian reductions and exact relation certificates}

\paragraph{1. Prove or refute the frozen $S_6$ and $S_8$ identities.}
The immediate arithmetic problem is to produce exact symbolic
relations for \eqref{cont26:eq:s6-recalled} and \eqref{cont26:eq:s8-candidate}.
The source's separating functional shows why repeating its
restricted single-product shuffle calculation cannot suffice for
$S_6$. A useful approach is to transport the convergent octahedral
relations used at weight five to weights seven and nine, and then
perform exact elimination. Each analytic relation schema should
include its branch, endpoint convergence, and regularization
justification before its matrix row is used. The desired output is a
short rational linear certificate whose target is the already frozen
vector, not another numerical relation search.

\paragraph{2. Find a general even-index reduction mechanism.}
Does \eqref{cont26:eq:even-index-question} hold for every $m\ge2$?
More specifically, can one derive the coefficient $-2$ of
$\beta(2m)\log2$ before doing weight-specific elimination?
An integral transformation that isolates that boundary term could
separate the universal part of the problem from the remaining
Gaussian double values. A recurrence in $m$, or a generating
function for the entire family, would be stronger than checking
additional individual weights. Conversely, an exact obstruction at
one weight would identify which new class of periods is required.

\paragraph{3. Compress and certify the relation search.}
The goal is not just to obtain a matrix of larger rank. Can the
octahedral, shuffle, stuffle, conjugation, and distribution relations
be organized into a small family of reusable certificates?
A valuable deliverable would be a verifier whose trusted analytic
inputs are explicitly listed and whose remaining work consists only
of rational arithmetic on words. Separating functionals should be
retained as diagnostics of missing relation families. Quotient
dimensions of a chosen presentation must continue to be
distinguished from dimensions of spaces of evaluated periods.

\subsection{Uniform asymptotics beyond the present error bounds}

\paragraph{4. Make the harmonic uniform constants effective.}
The proof of Theorem~\ref{cont26:thm:uniform} establishes universal
$C_J,p_J$, but does not tabulate useful explicit values. Derive
computable majorants for the derivatives of $H$, choose an explicit
localization radius, and propagate these bounds through the
Gaussian expansion. The practical target is a relative enclosure
for $R_{p,h}(a)$ from the saddle and a fixed number of derivatives,
valid without checking which ratio regime applies. Such a result
would complement the exact Euler method when the defining
alternating sum suffers severe cancellation.

\paragraph{5. Connect the bounded-$p$ and growing-$p$ depth limits.}
The new theorem allows $p/h\to0$ only when $p\to\infty$.
The source has a separate fixed-$p$ large-depth expansion.
Can one construct a single approximation with a relative error
uniform for $p$ in $[p_0,M(h)]$, including bounded $p$?
The change from an approximately Gaussian gamma law to a
finite-shape gamma law suggests keeping the latter expectation
unexpanded until the final step. A successful statement should
specify its dependence on a shift $a$ that also varies.

\paragraph{6. Replace the reflected inequality certificate by a
conceptual argument.}
Theorem~\ref{cont26:joint:thm:Q} has a complete finite proof, but the positive
quantity $D(x)$ deserves a structural explanation.
Can it be expressed as an integral with manifestly positive
integrand, a covariance of monotone functions, or the derivative
of a known gamma-function comparison? Such an argument could
give usable lower bounds on the curvature, and might extend to
$\log\Gamma(x+\alpha)$ after subtracting an appropriate affine
normalization. The affine subtraction and domain must be specified
so that the logarithms defining the phase remain real.

\paragraph{7. Obtain one reflected-moment expansion across all
exponent ratios.}
The proportional saddle expansion and the critical-window
expectation expansion overlap only after a further analysis.
The endpoint saddle has size $e^{-1/\lambda}$, so a hypothesis
keeping the saddle away from the boundary is unavailable.
Find a change of variable that resolves this exponential motion,
then prove a uniform remainder while $n,m\to\infty$ with no
compact restriction on $m/n$. Reflection should recover the
opposite endpoint. The parameter to match is
$\tau=m e^{-n/(m+1)}$, with the distinction between $m$ and
$m+1$ retained at correction order.

\paragraph{8. Classify when the fixed-exponent leading scale remains
valid, and extend the transition to all orders.}
Determine necessary and sufficient conditions, in a stated joint
growth domain, for $M_{n,m}/L_{n,m}\to1$.
The condition $m/n\to0$ alone is false, and the critical-window
theorem identifies a nonunit limiting profile. It does not by
itself settle every sequence outside that window.
For bounded $n/(m+1)-\log m$, derive the general coefficient form
\[
 \frac{M_{n,m}}{L_{n,m}}
 \sim e^{d\tau}\sum_{j\ge0}\frac{P_j(t,\tau)}{m^j},
\]
including the degree structure of $P_j$, a finite coefficient
algorithm, and a uniform remainder at each order.
The exact gamma moments already used here provide a route to
those polynomials; weighted tail control must be kept in the proof.

\paragraph{9. Study late coefficients when the reflected exponent
grows.}
The source's inverse-gamma coefficient analysis and optimal
truncation theorems have a separate saddle and a fixed reflected
power. Determine how the nearest singularity and least-term
remainder change when that power grows with the coefficient index.
This is not an automatic consequence of the joint moment theorem:
the coefficient-extraction contour and the moment integral have
different large parameters and different error budgets.

\subsection{Lerch discriminants and real-zero geometry}

\paragraph{10. Certify the local transition type at indices three and
four.}
The article proves an interior multiple root but does not prove
that it is exactly double. At the displayed numerical candidates,
certify a solution of
$(D_{n,1}^{\rho},\partial_aD_{n,1}^{\rho})=(0,0)$ in a small
rectangle and show
\[
 \partial_a^2D_{n,1}^{\rho}\ne0,\qquad
 \partial_\rho D_{n,1}^{\rho}\ne0.
\]
These inequalities would establish a transverse local creation or
annihilation of two real roots. A separate exclusion argument is
needed to prove that this is the only transition in $0<\rho<1$.
The endpoint logarithmic law helps determine a suitable parameter
scale, but its simple-root formula cannot be used at the multiple
root itself.

\paragraph{11. Determine the zero-count diagram for arbitrary
indices and derivative orders.}
For $n>k$, the general unfolding theorem reduces the small-parameter
count to the sign of $P_{n,k}(-\log2)$, with its nonvanishing stated
as a hypothesis. Can one determine that sign by elementary
inequalities over substantial ranges of $n,k$?
Can the endpoint $n$-zero count be proved uniformly in larger
families without separate numerical point selection?
Together these results would give lower bounds on the number of
interior discriminant events. They would also quantify the
derivative-order threshold beyond which the source's eventual
saturation holds for all $\rho$.

\paragraph{12. Resolve the full endpoint expansion and its interaction
with zeros.}
Theorem~\ref{cont26:thm:bif-rho-one} isolates the leading
$\epsilon^kL^{n+1}$ term. Determine the complete polynomial in
$L=\log(1/\epsilon)$ at order $\epsilon^k$, and continue to
higher powers of $\epsilon$ with explicit dependence on $a$.
One expects the lower logarithmic coefficients to record the
shift and the differentiated logarithmic polynomial that the
leading coefficient forgets. A Mellin-transform derivation should
be compared with the elementary split-sum proof, with all contour
shifts and multiple-pole contributions justified.
The resulting expansion could give sharper trajectories of
simple zeros and a starting point for simultaneous limits in
which the Stieltjes index also grows.

The research programme above connects analytic structure, finite
exact proof objects, and high-precision discovery. Each question
specifies which part of that work remains open and what would
constitute a checkable advance.

